\section{From Sentences to Disjunctive Entailment}
\label{sec:sentences}

This section reduces the truth of a sentence to finitely many disjunctive
entailments between flat constraints. The reduction is the same for every
class of trees that contains $\Bot$ and $\Top$ and is closed under the
constructor; we fix such a class $\Cls$.

\subsection{Normal Forms}\label{sec:sent-nf}

The truth value of a Boolean combination is computed from the truth values of
its component sentences, and a $\Pi_1$ sentence $\forall\bar x.\,\varphi$ is
true exactly when the $\Sigma_1$ sentence $\exists\bar x.\,\neg\varphi$ is
false. It therefore suffices to decide $\Sigma_1$ sentences.

A \emph{clause} is a conjunction of inequalities $s\le t$ and negated
inequalities $\neg(s\le t)$ between terms. Pushing negations inward and
distributing conjunctions over disjunctions turns a quantifier-free formula
$\varphi$ into an equivalent disjunction of clauses, a disjunctive normal form
of $\varphi$. Hence $\exists\bar x.\,\varphi$ is true over $\Cls$ exactly
when some clause of the normal form is \emph{realizable} in $\Cls$: some
assignment of trees of $\Cls$ to $\bar x$ makes all its conjuncts true.

\subsection{Naming Subterms}\label{sec:sent-flat}

Let $K$ be a clause. We introduce a variable $v_t$ for every term $t$ that
occurs in $K$, as a side of an inequality or as a subterm of one, with
$v_x=x$ for a variable $x$. The \emph{flat constraint} $\Gamma_K$ consists of
\begin{itemize}
\item $v_t=\Bot$ for $t=\Bot$, and $v_t=\Top$ for $t=\Top$, if these terms
  occur;
\item $f(v_{t_1},\dots,v_{t_n})\le v_t$ and $v_t\le f(v_{t_1},\dots,v_{t_n})$
  for every occurring term $t=f(t_1,\dots,t_n)$;
\item $v_s\le v_t$ for every inequality $s\le t$ of $K$.
\end{itemize}
The \emph{queries} of $K$ are the pairs $(v_s,v_t)$ for the negated
inequalities $\neg(s\le t)$ of $K$.

\begin{lemma}[Clauses]\label{lem:clause}
A clause $K$ with queries $(x_1,y_1),\dots,(x_k,y_k)$ is realizable in $\Cls$
if and only if $\Gamma_K\not\entails_{\Cls}x_1\le y_1\vee\dots\vee x_k\le
y_k$.
\end{lemma}

\begin{proof}
Let $\rho$ be an assignment in $\Cls$ of the variables of $K$ that realizes
$K$. Extend it by giving $v_t$ the value of the term $t$ under $\rho$; this
value lies in $\Cls$ because $\Cls$ contains $\Bot$ and $\Top$ and is closed
under the constructor. The first two groups of literals of $\Gamma_K$ hold
because the order is reflexive, the third because $\rho$ satisfies the
inequalities of $K$, and every query fails because $\rho$ satisfies the
negated inequalities. So the extension is a countermodel. Conversely, let
$\rho$ be a countermodel in $\Cls$. Since the order is antisymmetric, the two
literals for $t=f(t_1,\dots,t_n)$ give
$\rho(v_t)=f(\rho(v_{t_1}),\dots,\rho(v_{t_n}))$, so by induction on terms
$\rho(v_t)$ is the value of $t$ under the restriction of $\rho$ to the
variables of $K$. This restriction then satisfies every inequality of $K$,
by the third group, and no inequality $s\le t$ that is negated in $K$,
because $\rho$ refutes the query $(v_s,v_t)$.
\end{proof}

\begin{example}\label{ex:clause}
The sentence $\forall x.\;x\le\Bot\vee f(\Top,\Bot)\le x$ is the negation of
$\exists x.\;\neg(x\le\Bot)\wedge\neg(f(\Top,\Bot)\le x)$, whose body is a
single clause. Its flat constraint over the variables $x,b,t,a$ is
\[
  \Gamma=\{\,b=\Bot,\quad t=\Top,\quad f(t,b)\le a,\quad a\le f(t,b)\,\},
\]
with the queries $(x,b)$ and $(a,x)$. By \cref{lem:clause} the sentence is
true over $\Cls$ if and only if $\Gamma\entails_{\Cls}x\le b\vee a\le x$.
\end{example}

\begin{proposition}[Reduction to disjunctive entailment]\label{prop:sentences}
Suppose disjunctive entailment between flat constraints over $\Cls$ is
decidable, uniformly in the number of queries. Then it is decidable whether a
Boolean combination of $\Sigma_1$ and $\Pi_1$ sentences is true over $\Cls$.
Moreover, if two classes give every disjunctive entailment the same truth
value, they give every such sentence the same truth value.
\end{proposition}

\begin{proof}
By \cref{sec:sent-nf} it suffices to treat a $\Sigma_1$ sentence, which is
true exactly when one of the finitely many clauses of a normal form of its
body is realizable. By \cref{lem:clause} this is the failure of a disjunctive
entailment that is computed from the clause. The second statement follows
because the reduction does not depend on the class.
\end{proof}

The number of queries is the number of negated inequalities of a clause, so
the reduction needs disjunctive entailment for every $k$, and a single
decision procedure that takes $k$ as part of its input.

\subsection{Disjunctive Entailment Does Not Reduce to Plain Entailment}
\label{sec:sent-nonreduction}

Consider function types (\cref{ex:arrows}) and the entailment of
\cref{ex:clause},
\begin{equation}\label{eq:arrow-disj}
  \Gamma\;\entails\;x\le b\;\vee\;a\le x
  \qquad\text{with }b=\Bot\text{ and }a=\Top\to\Bot\text{ in every solution.}
\end{equation}
It holds, since every tree is $\Bot$ or lies above the least function type.
But $\Gamma$ entails neither $x\le b$, refuted by $x=\Top\to\Bot$, nor $a\le
x$, refuted by $x=\Bot$. For a covariant constructor, $f(\Bot,\dots,\Bot)$
plays the role of $\Top\to\Bot$. So deciding the disjuncts one at a time is
not enough, and countermodels of the single disjuncts do not help: a
countermodel of the disjunction must refute every disjunct with one
solution, and the two refuting values of $x$ above do not combine into one.
Moving negated disjuncts into the constraint takes us outside plain
entailment as well: $\Gamma\entails\varphi_1\vee\varphi_2$ holds if and only
if $\Gamma\wedge\neg\varphi_1\entails\varphi_2$, but $\neg\varphi_1$ is not
a constraint, and plain entailment allows only constraints on the left.
The $k$ failures therefore have to be decided jointly. The next section
describes what it means for $k$ failures to be realized by one solution.
